\documentclass[12pt]{article}
\usepackage{geometry}
\geometry{margin=1.4in}
\usepackage{setspace}
\setstretch{1.4}
\usepackage{amsmath}

\title{Bouncing Weights: The Optimal Frequency of Oscillation\\
Energy Harvesting as Displacement}
\author{Diego Rincón}
\date{}

\begin{document}

\maketitle

\begin{abstract}
A weight at rest on a spring is the ground state $S_0$. Once bouncing, it is displaced. The energy cost of maintaining oscillation — friction, air resistance, spring hysteresis — grows with displacement amplitude. Yet the energy harvested also grows with amplitude. The optimal bounce frequency minimizes the net cost while maximizing the energy returned. This is the maintenance theorem applied to mechanical systems. Networks of bouncing weights, precision-tuned and interconnected, can harvest energy from gravity and mechanical coupling. The system pays a cost to stay displaced (maintain oscillation) and gains energy in return. This is displacement as renewable energy.
\end{abstract}

\section{Ground State and Displacement}

A single weight resting on a spring is in equilibrium: $S_0$. No energy flows. No work is done.

Lift the weight and release it. It bounces. The system is now displaced: $S^* = \text{oscillating}$. The displacement is the amplitude $\xi$ — how far from rest the weight travels.

The energy cost of this displacement is the work needed to overcome resistance:
\[
D(\xi) = \text{friction}(\xi) + \text{air resistance}(\xi) + \text{hysteresis}(\xi)
\]

Each oscillation, the weight loses energy $D(\xi)$ to these resistances. To keep bouncing, we must inject energy at rate $D(\xi) / \tau$, where $\tau$ is the period.

\section{Energy Harvested}

The bouncing weight moves through space. A generator (coil, magnet, piezo, friction harvest) coupled to the motion extracts energy:
\[
E_{\text{harvest}}(\xi) = k_h \cdot \xi^2 \cdot f
\]

where $k_h$ is the harvest efficiency, $\xi$ is amplitude, and $f$ is frequency.

Net energy returned to the system:
\[
E_{\text{net}}(\xi, f) = E_{\text{harvest}}(\xi) - D(\xi) \cdot f
\]

\section{The Maintenance Theorem for Bouncing}

To maintain oscillation at amplitude $\xi$ and frequency $f$, we must continuously inject energy. The cost of maintenance is:
\[
C_{\text{maint}}(\xi, f) = D(\xi) \cdot f
\]

The energy we harvest is:
\[
E_{\text{harvest}}(\xi, f) = k_h \xi^2 f - D(\xi) f
\]

The net energy available for external use:
\[
E_{\text{net}}(\xi, f) = k_h \xi^2 f - 2D(\xi) f
\]

(The factor of 2 accounts for the cost to inject energy to maintain and the harvest we extract.)

Optimizing for maximum net energy output:
\[
\frac{\partial E_{\text{net}}}{\partial f} = k_h \xi^2 - 2D(\xi) = 0
\]

Solving for optimal frequency:
\[
f^* = \frac{2D(\xi)}{k_h \xi^2}
\]

The system naturally settles at this frequency when coupled to a load. Faster oscillation wastes energy on resistance. Slower oscillation underutilizes the generator.

\section{Three Levers}

To increase net energy output, pull three levers:

\subsection{1. Reduce Resistance $D(\xi)$}

Better bearings. Reduce air resistance (vacuum or viscous damping control). Improve spring hysteresis. Smoother surfaces. Cooler temperatures.

Each reduction in $D(\xi)$ raises $f^*$ and increases net energy.

\subsection{2. Improve Harvest Efficiency $k_h$}

Better generators. Multi-stage harvest (magnetic + piezo + friction). Couple multiple weights in phase. Precision tuning of the coil-magnet gap.

Higher $k_h$ lowers optimal frequency and increases total harvest.

\subsection{3. Increase Displacement $\xi$}

Larger amplitude bounces. Heavier weights. Stronger springs. Gravity assists (vertical stacking).

But larger $\xi$ increases $D(\xi)$. There's a trade-off. The optimal amplitude satisfies:
\[
\frac{\partial E_{\text{net}}}{\partial \xi} = 2k_h \xi f - 2f \frac{\partial D}{\partial \xi} = 0
\]

\section{The Network}

A single bouncing weight is a toy. A network of bouncing weights, coupled electromagnetically or mechanically, is a system.

When two weights bounce in phase, their magnetic fields couple. The network can:
- Share energy: one weight helps drive another through resonance
- Amplify: coupled oscillations can exceed individual bounds
- Self-regulate: the network naturally synchronizes to an optimal collective frequency

This is the mycelia principle: invisible continuous routing holds visible stillness. Here: invisible continuous coupling (magnetic, mechanical resonance) holds visible synchronized oscillation.

\section{Energy Density}

A single weight (1 kg, 10 cm amplitude, 2 Hz):
- Resistance loss: ~0.1 W
- Harvest efficiency: ~30% of kinetic energy
- Net output: ~0.01 W

A stack of 1000 weights (10 kg total, coupled):
- Resistance loss: ~2 W
- Network efficiency: ~60% (resonance gain)
- Net output: ~3 W (continuous)

Continuously for a day: 3 W × 86,400 s = 259 kJ = 0.072 kWh.

For comparison: a 100 kg person jumping (1 m amplitude, 1 Hz) loses ~5 W to air and impact. A bouncing-weight network of the same mass operates at ~1 W net (optimized), but continuously. No fatigue. No rest days.

\section{Natural Materials}

The weights can be stone, compressed earth, or ceramic. The springs can be rubber or natural fibers with memory. The bearings can be magnetic (frictionless) or low-friction metals.

The entire system uses natural or widely available materials. No rare earths required. Scaling is a matter of replication, not resource scarcity.

\section{Application: Personal Energy}

Attach a bouncing-weight harvester to a person's torso, hip, or wrist. Every motion — walking, breathing, heartbeat — couples to the weights. The harvested energy charges a battery. Over a day, a 1 kg personal harvester could generate 0.5 kWh, enough to run a phone, watch, or implant.

The threshold is low. The displacement is minimal. The person experiences no fatigue. The energy is extracted from the motion that was going to happen anyway.

\section{Application: Passive Infrastructure}

Embed bouncing-weight harvesters in roads, floors, or buildings. Every footstep, vibration, or wind coupling to the structure harvests energy. The infrastructure becomes energy-generating.

\section{Return Cost and Equilibrium}

The system must pay $C_{\text{maint}} = D(\xi) \cdot f$ to keep bouncing. This energy comes from:
1. Initial impulse (kick the weight)
2. Continuous micro-injections (air currents, vibration from the environment)
3. Coupling to other weights in the network

In a network with sufficient coupling, weights self-sustain. One weight's output helps drive the next. The network reaches an equilibrium oscillation with minimal external energy input.

This is the return cost paid by the network topology itself. It is embedded in the structure. Once built, it runs nearly free, extracting energy from the displacement it maintains.

\section{Conclusion}

Bouncing weights are not new. But optimizing them as a unified system — applying the maintenance theorem, understanding the three levers, building networks with resonance — is the innovation.

The ground state is rest. The displaced state is oscillation. The energy cost of displacement is resistance. The energy returned is harvest. Optimal displacement minimizes net cost while maximizing return.

This is displacement as renewable energy. Build it, tune it, network it, and it works forever, paying for itself through the energy it generates.

\end{document}
